%% maestro2tex -- generated figure; do not edit by hand.
\providecommand{\MaestroRoot}{include/analog/}
\begin{figure}[htbp]
  \centering
  {\pgfplotsset{compat=1.18}%
  \begin{tikzpicture}
    \begin{axis}[
      width=0.82\linewidth, height=6.2cm,
      xlabel={Frequency~[\si{\hertz}]}, ylabel={Loop phase $\ang{180}+\angle A$~[\si{\degree}]},
      grid=both,
      major grid style={line width=0.2pt, draw=black!14},
      minor grid style={line width=0.2pt, draw=black!7},
      minor tick num=1,
      axis line style={line width=0.4pt, draw=black!45},
      tick style={line width=0.4pt, draw=black!45},
      tick align=outside,
      label style={font=\small}, tick label style={font=\small},
      enlarge x limits=0.02, enlarge y limits=0.10,
      every axis plot/.append style={line join=round, no marks},
      legend style={font=\footnotesize, draw=black!25, fill=white,
                    fill opacity=0.9, text opacity=1, row sep=1pt},
      legend cell align=left, legend pos=south west,
      xmode=log,
      scaled y ticks=false,
      y tick label style={/pgf/number format/fixed,
                          /pgf/number format/precision=0},
      ymin=0, ymax=200, xmax=1e8
    ]
      \addplot[mtxA, line width=1.1pt, solid] table {\MaestroRoot data/tia_loopphase_00.dat};
      \addlegendentry{tt}
      \addplot[mtxB, line width=1.1pt, dashed] table {\MaestroRoot data/tia_loopphase_01.dat};
      \addlegendentry{ff}
      \addplot[mtxC, line width=1.1pt, dotted] table {\MaestroRoot data/tia_loopphase_02.dat};
      \addlegendentry{fs}
      \addplot[mtxD, line width=1.1pt, dashdotted] table {\MaestroRoot data/tia_loopphase_03.dat};
      \addlegendentry{sf}
      \addplot[mtxE, line width=1.1pt, densely dashed] table {\MaestroRoot data/tia_loopphase_04.dat};
      \addlegendentry{ss}
    \end{axis}
  \end{tikzpicture}}
  \caption[{Phase of the same loop-gain measurement, over the five process corners.}]{Phase of the same loop-gain measurement, over the five process corners. Spectre's \texttt{stb} analysis reports $\ang{180}+\angle A$, so the curve starts at \ang{180} at DC and its value at the \SI{6.4}{\mega\hertz} crossover \emph{is} the phase margin. The dominant pole costs \ang{90} by \SI{100}{\kilo\hertz}; the lift from \ang{90.3} there back to \ang{92.3} at \SI{1}{\mega\hertz} is the left-half-plane zero of Equation~(\ref{eq:tia-zero}) adding lead, which is what holds the margin near \ang{88} through crossover despite the second pole at \SI{10.8}{\mega\hertz}.}
  \label{fig:tia-loopphase}
\end{figure}
